\documentclass[11pt,a4paper]{article}
\usepackage[utf8]{inputenc}
\usepackage[spanish]{babel}
\usepackage{amsmath, amssymb}
\usepackage{geometry}
\geometry{top=2.5cm, bottom=2.5cm, left=2.5cm, right=2.5cm}

\begin{document}

\section*{Piense, dialogue y comparta}

\begin{itemize}
    \item[\textbf{1.}] Usted es miembro de un grupo de testigos expertos que brinda servicios científicos a la comunidad legal. Se ha pedido a su grupo que defienda a un equipo de futbol de la NFL, que ha quedado atrapado en una práctica embarazosa: han estado llenando los balones de sus equipos con helio en lugar de aire, creyendo que el helio proporcionaría una mayor fuerza de empuje en los balones, causando que sus pases y patadas sean más largos. A pesar de su descontento con los esfuerzos del equipo para obtener una ventaja injusta, su grupo legal cree que todos merecen una defensa, por lo que acepta llevar el caso.
    \begin{enumerate}
        \item[\textbf{(a)}] Desarrolle un argumento para sostener que llenar los balones de futbol con helio no proporcionaría una fuerza adicional de flotación en los balones; por tanto, el equipo no estaba tratando de obtener una ventaja.
        \item[\textbf{(b)}] Desarrolle un argumento privado para presentarle al equipo sosteniendo que llenar las bolas con helio realmente disminuirá el rendimiento.
    \end{enumerate}
    \vspace*{9cm}

    \newpage

    \item[\textbf{2.}] Es un día cálido, y un estudiante decide que le gustaría pasar unas horas nadando en la piscina. Usando su equipo de snorkel, él ve el fondo de la piscina. Después de un rato, observa un tubo de PVC apoyado contra la casa y piensa en usarlo como un tubo largo. Él saca la boquilla de su tubo y la sujeta al extremo de la tubería de PVC. ¡Su plan es tener la boca en la boquilla del tubo largo y hundirse en el extremo profundo de la piscina, respirando hasta el fondo! Mientras prepara su largo snorkel, llega su tío, que es neumólogo, para una visita. Él le pregunta al estudiante qué está haciendo con el tubo y él explica. El tío está horrorizado y dice que es muy peligroso intentar esta actividad. El estudiante está decepcionado, pero luego le dice a su tío que aguantará la respiración mientras se hunde hasta el fondo, manteniendo la boquilla cerrada con el pulgar, coloca la boquilla mientras llega al final y comienza a respirar. ¡Su tío parece aún más horrorizado, y le dice que esa actividad podría ser fatal! Discuta en su grupo: ¿Por qué el buceo profundo es tan peligroso?
    \vspace*{9cm}

    \newpage

    \item[\textbf{3.}] \textbf{ACTIVIDAD} Realice las siguientes actividades y analice la física detrás de los resultados con los miembros de su grupo: 
    \begin{enumerate}
        \item Coloque una lata de refresco dietético y refresco común de la misma marca en un recipiente grande con agua. ¿Ambos flotan? Explique los resultados. 
        \item Abra una lata de bebida carbonatada transparente y llene un vaso transparente con el líquido. Ahora vierta unas pocas pasas en la bebida y registre su comportamiento.
    \end{enumerate}
    \vspace*{8cm}
\end{itemize}

\newpage

\section*{Problemas}

\subsection*{SECCIÓN 10.1 Presión}

\begin{itemize}
    \item[\textbf{1.}] Un hombre grande se sienta en una silla de cuatro patas, con sus pies levantados del piso. La masa combinada del hombre y la silla es \(95.0\,\text{kg}\). Si las patas de la silla son circulares y tienen un radio de \(0.500\,\text{cm}\) en el fondo, ¿qué presión ejerce cada pata sobre el piso?
    \vspace*{6cm}
    
    \item[\textbf{2.}] El núcleo de un átomo se puede modelar como constituido por protones y neutrones cercanamente empaquetados. Cada partícula tiene una masa de \(1.67 \times 10^{-27}\,\text{kg}\) y un radio del orden de \(10^{-15}\,\text{m}\).
    \begin{enumerate}
        \item Utilice este modelo y los datos proporcionados para estimar la densidad del núcleo de un átomo.
        \item Compare su resultado con la densidad de un material, por ejemplo, el hierro. ¿Qué sugieren su resultado y la comparación respecto a la estructura de la materia?
    \end{enumerate}
    \vspace*{7cm}
    
    \item[\textbf{3.}] Estime la masa total de la atmósfera de la Tierra. (El radio de la Tierra es \(6.37 \times 10^6\,\text{m}\), y la presión atmosférica en la superficie es \(1.013 \times 10^5\,\text{Pa}\).)
    \vspace*{6cm}
\end{itemize}

\newpage

\subsection*{SECCIÓN 10.2 Variación de la presión con la profundidad}

\begin{itemize}
    \item[\textbf{4.}] ¿Por qué no es posible la siguiente situación? La figura P10.4 muestra a Superman intentando beber agua fría con un popote de longitud \(\ell = 12.0\,\text{m}\). Las paredes del popote tubular son muy fuertes y no colapsan. Con su gran fuerza, logra la máxima succión posible y disfruta beber el agua fría.
    
    % [Espacio para Figura P10.4]
    \vspace*{6cm}
    
    \item[\textbf{5.}] ¿Cuál debe ser el área de contacto entre una ventosa (completamente vacía) y un techo, si la ventosa debe soportar el peso de un estudiante de \(80.0\,\text{kg}\)?
    \vspace*{6cm}
    
    \item[\textbf{6.}] Para el sótano de una nueva casa, se cava un hoyo en el suelo, con lados verticales que bajan \(2.40\,\text{m}\). Una pared de cimiento de concreto se construye horizontal los \(9.60\,\text{m}\) de ancho de la excavación. Esta pared de cimiento mide \(0.183\,\text{m}\) desde el frente del hoyo del sótano. Durante una tormenta, el drenaje de la calle llena el espacio enfrente de la pared de concreto, pero no el sótano detrás de la pared. El agua no se filtra en el suelo de arcilla. Encuentre la fuerza que ejerce el agua sobre la pared de cimiento. En comparación, el peso del agua está dado por \(2.40\,\text{m} \times 9.60\,\text{m} \times 0.183\,\text{m} \times 1000\,\text{kg/m}^3 \times 9.80\,\text{m/s}^2 = 41.3\,\text{kN}\).
    \vspace*{7cm}
    
    \item[\textbf{7.}] \textbf{Problema de repaso.} Una esfera sólida de latón (módulo volumétrico de \(10.0 \times 10^{10}\,\text{N/m}^2\)) con un diámetro de \(3.00\,\text{m}\) es lanzada al océano. ¿Cuánto disminuye el diámetro de la esfera cuando se sumerge a una profundidad de \(1.00\,\text{km}\)?
    \vspace*{7cm}
\end{itemize}

\newpage

\subsection*{SECCIÓN 10.3 Mediciones de presión}

\begin{itemize}
    \item[\textbf{8.}] El cerebro humano y la médula espinal están sumergidos en el fluido cerebroespinal. El fluido normalmente es continuo entre las cavidades craneal y espinal y ejerce una presión de \(100\) a \(200\,\text{mm}\) de \(\text{H}_2\text{O}\) sobre la presión atmosférica prevaleciente. En el trabajo médico, las presiones usualmente se miden en milímetros de \(\text{H}_2\text{O}\) porque los fluidos corporales, incluido el fluido cerebroespinal, por lo general tienen la misma densidad que el agua. La presión del fluido cerebroespinal se puede medir mediante una sonda espinal, como se ilustra en la figura P10.8. Un tubo hueco se inserta en la columna vertebral y se observa la altura a la que se eleva el fluido. Si el fluido se eleva a una altura de \(160\,\text{mm}\), su presión manométrica se escribe como \(160\,\text{mm H}_2\text{O}\).
    \begin{enumerate}
        \item Exprese esta presión en pascales, en atmósferas y en milímetros de mercurio.
        \item Algunas condiciones que bloquean o inhiben el flujo de fluido cerebroespinal se pueden investigar mediante la prueba de Queckenstedt. En este procedimiento, se comprimen las venas en la nuca del paciente para hacer que la presión sanguínea se eleve en el cerebro, lo cual se transmite al fluido cerebroespinal. Explique cómo el nivel de fluido en la sonda espinal puede utilizarse como herramienta de diagnóstico para la condición de la espina del paciente.
    \end{enumerate}
    
    % [Espacio para Figura P10.8]
    \vspace*{8cm}
    
    \item[\textbf{9.}] Blaise Pascal duplicó el barómetro de Torricelli usando un vino tinto Bordeaux, de \(984\,\text{kg/m}^3\) de densidad, como el líquido de trabajo (figura P10.9).
    \begin{enumerate}
        \item ¿Cuál fue la altura \(h\) de vino para presión atmosférica normal?
        \item ¿Esperaría que el vacío sobre la columna sea tan bueno como para el mercurio?
    \end{enumerate}
    
    % [Espacio para Figura P10.9]
    \vspace*{7cm}
    
    \item[\textbf{10.}] Un tanque con un fondo plano de área \(A\) y lados verticales se llena con agua con una profundidad \(h\). La presión es \(P_0\) en la superficie superior.
    \begin{enumerate}
        \item ¿Cuál es la presión absoluta en el fondo del tanque?
        \item Suponga que un objeto de masa \(M\) y densidad menor a la densidad del agua se coloca en el tanque y flota. No se desborda agua. ¿Cuál es el aumento resultante de presión en el fondo del tanque?
    \end{enumerate}
    \vspace*{7cm}
\end{itemize}

\newpage

\subsection*{SECCIÓN 10.4 Fuerzas de flotación y principio de Arquímedes}

\begin{itemize}
    \item[\textbf{11.}] La fuerza gravitacional que se ejerce sobre un objeto sólido es \(5.00\,\text{N}\). Cuando el objeto se suspende de una balanza de resorte y se sumerge en agua, la lectura en la balanza es \(3.50\,\text{N}\) (figura P10.11). Encuentre la densidad del objeto.
    
    % [Espacio para Figura P10.11]
    \vspace*{6cm}
    
    \item[\textbf{12.}] Un bloque metálico de \(10.0\,\text{kg}\) que mide \(12.0\,\text{cm}\) por \(10.0\,\text{cm}\) por \(10.0\,\text{cm}\), está suspendido de una balanza y sumergido en agua, como se muestra en la figura P10.11b. La dimensión de \(12.0\,\text{cm}\) es vertical y la parte superior del bloque está \(5.00\,\text{cm}\) debajo de la superficie del agua.
    \begin{enumerate}
        \item ¿Cuáles son las magnitudes de las fuerzas que actúan sobre las partes superior e inferior del bloque debidas al agua circundante?
        \item ¿Cuál es la lectura en la balanza de resorte?
        \item Demuestre que la fuerza de flotación es igual a la diferencia entre las fuerzas sobre las partes inferior y superior del bloque.
    \end{enumerate}
    \vspace*{8cm}
    
    \item[\textbf{13.}] Una esfera plástica flota en agua con \(50.0\%\) de su volumen sumergido. Esta misma esfera flota en glicerina con \(40.0\%\) de su volumen sumergido. Determine las densidades de \textbf{(a)} la glicerina y \textbf{(b)} la esfera.
    \vspace*{6cm}
\end{itemize}

\newpage

\begin{itemize}
    \item[\textbf{14.}] El peso de un bloque rectangular de material de baja densidad es \(15.0\,\text{N}\). Con una cuerda delgada, el centro de la cara inferior horizontal del bloque se amarra al fondo de un vaso de precipitados parcialmente lleno con agua. Cuando \(25.0\%\) del volumen del bloque está sumergido, la tensión en la cuerda es \(10.0\,\text{N}\).
    \begin{enumerate}
        \item Encuentre la fuerza de flotación sobre el bloque.
        \item Ahora al vaso de precipitados se le agrega sin interrupción aceite de \(800\,\text{kg/m}^3\) de densidad, lo que forma una capa sobre el agua y rodea al bloque. El aceite ejerce fuerzas sobre cada una de las cuatro paredes laterales del bloque que el aceite toca. ¿Cuáles son las direcciones de dichas fuerzas?
        \item ¿Qué le ocurre a la tensión en la cuerda conforme se agrega el aceite? Explique cómo el aceite tiene este efecto sobre la tensión de la cuerda.
        \item La cuerda se rompe cuando su tensión alcanza \(60.0\,\text{N}\). En este momento, \(25.0\%\) del volumen del bloque todavía está bajo la línea del agua. ¿Qué fracción adicional del volumen del bloque continúa por debajo de la superficie superior del aceite?
    \end{enumerate}
    \vspace*{9cm}
    
    \item[\textbf{15.}] Un bloque de madera de volumen \(5.24 \times 10^{-4}\,\text{m}^3\) flota en agua, y un pequeño objeto de acero de masa \(m\) se coloca en la parte superior del bloque. Cuando \(m = 0.310\,\text{kg}\), el sistema está en equilibrio y la parte superior del bloque de madera está justo al nivel del agua.
    \begin{enumerate}
        \item ¿Cuál es la densidad de la madera?
        \item ¿Qué le ocurre al bloque cuando el objeto de acero es reemplazado por un objeto cuya masa es menor que \(0.310\,\text{kg}\)?
        \item ¿Qué le sucede al bloque cuando el objeto de acero se sustituye por un objeto cuya masa es mayor que \(0.310\,\text{kg}\)?
    \end{enumerate}
    \vspace*{7cm}
    
    \item[\textbf{16.}] Un hidrómetro es un instrumento utilizado para determinar la densidad de los líquidos. En la figura P10.16 se muestra uno simple. El bulbo de una jeringa se presiona y libera para dejar que la atmósfera eleve una muestra del líquido de interés en un tubo que contiene una barra calibrada de densidad conocida. La barra, de longitud \(L\) y densidad promedio \(\rho_0\), flota parcialmente sumergida en el líquido de densidad \(\rho\). Una longitud \(h\) de la barra sobresale de la superficie del líquido. Demuestre que la densidad del líquido está dada por
    \[
    \rho = \frac{\rho_0 L}{L - h}
    \]
    
    % [Espacio para Figura P10.16]
    \vspace*{7cm}
\end{itemize}

\newpage

\begin{itemize}
    \item[\textbf{17.}] Remítase al problema 16 y la figura P10.16. Se construirá un hidrómetro con una barra cilíndrica flotante. Se colocarán nueve marcas a lo largo de la barra para indicar densidades que tengan valores de \(0.98\,\text{g/cm}^3\), \(1.00\,\text{g/cm}^3\), \(1.02\,\text{g/cm}^3\), \(1.04\,\text{g/cm}^3\), \dots, \(1.14\,\text{g/cm}^3\). La hilera de marcas comenzará \(0.200\,\text{cm}\) desde el extremo superior de la barra y terminará \(1.80\,\text{cm}\) desde el extremo superior.
    \begin{enumerate}
        \item ¿Cuál es la longitud requerida de la barra?
        \item ¿Cuál debe ser su densidad promedio?
        \item ¿Las marcas deberían estar igualmente espaciadas? Explique su respuesta.
    \end{enumerate}
    \vspace*{8cm}
    
    \item[\textbf{18.}] El 21 de octubre de 2001, Ian Ashpole del Reino Unido logró un récord de altura de \(3.35\,\text{km}\) (\(11000\,\text{pies}\)) impulsado por 600 globos llenos con helio. Cada globo lleno tuvo un radio de aproximadamente \(0.50\,\text{m}\) y una masa estimada de \(0.30\,\text{kg}\).
    \begin{enumerate}
        \item Estime la fuerza de flotación total sobre los 600 globos.
        \item Estime la fuerza neta hacia arriba sobre los 600 globos.
        \item Ashpole utilizó paracaídas para su retorno a la Tierra después de que los globos empezaron a reventarse a elevadas altitudes y disminuyera la fuerza boyante. ¿Por qué se reventaron los globos?
    \end{enumerate}
    \vspace*{7cm}
    
    \item[\textbf{19.}] Usted tiene un trabajo en una empresa que produce suministros para fiestas. Está diseñando globos llenos de helio para vender como regalos. Para ahorrar dinero en costos de producción, parte de este diseño es elegir de entre una selección disponible la ficha menos masiva necesaria para ser atada al extremo inferior de la cuerda que cuelga del globo para evitar que éste se levante de una mesa, bandeja de comida de hospital, tocador de dormitorio, etc. De esta manera, la ficha permanece estacionaria en la superficie plana y el globo se mantiene flotando por encima de la ficha a una altura fija, con la cuerda recta. Está trabajando en un globo cuya envoltura es muy delgada y tiene una masa de \(0.150\,\text{kg}\). La envoltura se llena a un volumen de \(0.250\,\text{m}^3\) con helio a presión atmosférica. La cuerda tiene una masa de \(0.0700\,\text{kg}\). Entre la selección de fichas se encuentran aquellas con masas de \(10.0\,\text{g}\), \(20.0\,\text{g}\), \(30.0\,\text{g}\), \(40.0\,\text{g}\) y \(50.0\,\text{g}\). Elija la ficha apropiada para el globo en el que está trabajando.
    \vspace*{7cm}
\end{itemize}

\newpage

\subsection*{SECCIÓN 10.5 Dinámica de fluidos}

\begin{itemize}
    \item[\textbf{20.}] El agua que fluye de una manguera de jardín de diámetro \(2.74\,\text{cm}\) llena una cubeta de \(25\,\text{L}\) en \(1.50\,\text{min}\).
    \begin{enumerate}
        \item ¿Cuál es la rapidez del agua que sale del extremo de la manguera?
        \item Se acopla una boquilla al extremo de la manguera. Si el diámetro de la boquilla es un tercio del diámetro de la manguera, ¿cuál es la rapidez del agua que sale por la boquilla?
    \end{enumerate}
    \vspace*{7cm}
    
    \item[\textbf{21.}] Sobre un dique de altura \(h\) cae agua con un caudal de masa \(R\), en unidades de kilogramos por segundo.
    \begin{enumerate}
        \item Demuestre que la potencia disponible a causa del agua es 
        \[
        P = Rgh
        \]
        donde \(g\) es la aceleración en caída libre.
        \item Cada unidad hidroeléctrica en el dique Grand Coulee toma agua en una cantidad de \(8.50 \times 10^5\,\text{kg/s}\) desde una altura de \(87.0\,\text{m}\). La potencia desarrollada por la caída de agua se convierte en energía eléctrica con una eficiencia de \(85.0\%\). ¿Cuánta energía eléctrica produce cada unidad hidroeléctrica?
    \end{enumerate}
    \vspace*{7cm}
\end{itemize}

\subsection*{SECCIÓN 10.6 Ecuación de Bernoulli}

\begin{itemize}
    \item[\textbf{22.}] Un legendario niño holandés salvó a Holanda al poner su dedo en un hoyo de \(1.20\,\text{cm}\) de diámetro en un dique. Si el hoyo estaba \(2.00\,\text{m}\) bajo la superficie del Mar del Norte (densidad \(1030\,\text{kg/m}^3\)),
    \begin{enumerate}
        \item ¿cuál fue la fuerza sobre su dedo?
        \item Si él hubiera sacado el dedo del hoyo, ¿durante qué intervalo de tiempo, el agua liberada llenaría 1 acre de tierra a una profundidad de 1 pie? Suponga que el hoyo mantuvo constante su tamaño.
    \end{enumerate}
    \vspace*{7cm}
\end{itemize}

\newpage

\begin{itemize}
    \item[\textbf{23.}] Desde el Río Colorado se bombea agua para suministrar a Grand Canyon Village, ubicada a la orilla del cañón. El río está a una elevación de \(564\,\text{m}\) y la villa está a una elevación de \(2096\,\text{m}\). Imagine que el agua se bombea a través de una larga tubería de \(15.0\,\text{cm}\) de diámetro, impulsada por una bomba en el extremo inferior.
    \begin{enumerate}
        \item ¿Cuál es la presión mínima a la que el agua debe bombearse si ha de llegar a la villa?
        \item Si \(4500\,\text{m}^3\) de agua se bombean por día, ¿cuál es la rapidez del agua en la tubería?
    \end{enumerate}
    \textit{Nota:} Suponga que la aceleración en caída libre y la densidad del aire son constantes en este intervalo de elevaciones. Las presiones que calcule son muy altas para una tubería ordinaria. En realidad el agua se eleva en etapas mediante varias bombas a través de tuberías cortas.
    \vspace*{7cm}
    
    \item[\textbf{24.}] En flujo ideal, un líquido de \(850\,\text{kg/m}^3\) de densidad se mueve desde un tubo horizontal de \(1.00\,\text{cm}\) de radio a un segundo tubo horizontal de \(0.500\,\text{cm}\) de radio a la misma elevación que el primer tubo. La presión difiere por \(\Delta P\) entre el líquido en un tubo y el líquido en el segundo tubo.
    \begin{enumerate}
        \item Encuentre la rapidez de flujo volumétrico como una función de \(\Delta P\).
        \item Evalúe la rapidez de flujo volumétrico para \(\Delta P = 6.00\,\text{kPa}\) y \textbf{(c)} \(\Delta P = 12.0\,\text{kPa}\).
    \end{enumerate}
    \vspace*{7cm}
    
    \item[\textbf{25.}] \textbf{Problema de repaso.} El géiser Old Faithful en el parque nacional Yellowstone erupciona a intervalos aproximados de 1 hora, y la altura de la columna de agua alcanza \(40.0\,\text{m}\) (figura P10.25).
    \begin{enumerate}
        \item Modele la corriente que se eleva como una serie de gotas separadas. Analice el movimiento en caída libre de una de las gotas para determinar la rapidez a la que el agua deja el suelo.
        \item \textbf{¿Qué pasaría si?} Modele la corriente que se eleva como un fluido ideal en un flujo de líneas de corriente. Use la ecuación de Bernoulli para determinar la rapidez del agua mientras deja el nivel del suelo.
        \item ¿De qué modo se compara la respuesta del inciso (a) con la respuesta del inciso (b)?
        \item ¿Cuál es la presión (arriba de la atmosférica) en la cámara subterránea caliente si su profundidad es de \(175\,\text{m}\)? Suponga que la cámara es grande en comparación con la boca del géiser.
    \end{enumerate}
    
    % [Espacio para Figura P10.25]
    \vspace*{8cm}
\end{itemize}

\newpage

\begin{itemize}
    \item[\textbf{26.}] Está trabajando como testigo experto para el propietario de un complejo de rascacielos en el centro de la ciudad. El propietario está siendo demandado por peatones en las calles debajo de sus edificios que resultaron heridos por la caída de cristales cuando ventanas salieron expelidas de los lados del edificio. El efecto Bernoulli puede tener consecuencias importantes para las ventanas en dichos edificios. Por ejemplo, el viento puede soplar alrededor de un rascacielos a una velocidad notablemente alta, creando una baja presión en la superficie exterior de las ventanas. La presión atmosférica más alta en el aire inmóvil dentro de los edificios puede hacer que las ventanas sean expulsadas.
    \begin{enumerate}
        \item En su investigación sobre el caso, encontrará algunas vistas generales del proyecto de su cliente, como se muestra a continuación. El proyecto incluye dos altos rascacielos y un área de parque en un terreno cuadrado. El plan (i) (figura P10.26, i) fue presentado por los arquitectos y planificadores originales. En el último minuto, el propietario decidió que no quería que los terrenos del parque se dividieran en dos áreas y presentó el Plan (ii) (figura P10.26, ii), que es la forma en que se construyó el proyecto. Explique a su cliente por qué el plan (ii) es una situación mucho más peligrosa en términos de ventanas expulsadas que el plan (i).
        \item Su cliente no está convencido por su argumento conceptual en el inciso (a), por lo que proporciona un argumento numérico. Un viento horizontal sopla con una velocidad de \(11.2\,\text{m/s}\) fuera de un gran panel de vidrio con dimensiones de \(4.00\,\text{m} \times 1.50\,\text{m}\). Suponga que la densidad del aire es constante a \(1.20\,\text{kg/m}^3\). El aire dentro del edificio está a presión atmosférica. Calcule la fuerza total ejercida por el aire sobre el cristal de la ventana para su cliente.
        \item \textbf{¿Qué pasaría si?} Para convencer aún más a su cliente de los problemas con el diseño del edificio, calcule la fuerza total ejercida por el aire sobre el cristal si la velocidad del viento los edificios son \(22.4\,\text{m/s}\), dos veces más altos que en el inciso (b).
    \end{enumerate}
    
    % [Espacio para Figura P10.26]
    \vspace*{9cm}
\end{itemize}

\subsection*{SECCIÓN 10.7 Flujo de fluidos viscosos en tuberías}

\begin{itemize}
    \item[\textbf{27.}] Se ha colocado un recubrimiento fino de glicerina de \(1.50\,\text{mm}\) entre dos portaobjetos de microscopio de \(1.00\,\text{cm}\) de ancho y \(4.00\,\text{cm}\) de largo. Encuentre la fuerza requerida para jalar de uno de los portaobjetos del microscopio a una velocidad constante de \(0.300\,\text{m/s}\) en relación con el otro portaobjetos.
    \vspace*{6cm}
    
    \item[\textbf{28.}] Una aguja hipodérmica mide \(3.00\,\text{cm}\) de largo y \(0.300\,\text{mm}\) de diámetro. ¿Qué diferencia de presión entre la entrada y la salida de la aguja es necesaria para que el caudal de agua a través de ella sea de \(1.00\,\text{g/s}\)? (Use \(1.00 \times 10^{-3}\,\text{Pa}\cdot\text{s}\) como la viscosidad del agua.)
    \vspace*{6cm}
\end{itemize}

\newpage

\begin{itemize}
    \item[\textbf{29.}] ¿Qué radio de aguja se debe usar para inyectar un volumen de \(500\,\text{cm}^3\) de una solución en un paciente en \(30.0\,\text{minutos}\)? Suponga que la longitud de la aguja es de \(2.50\,\text{cm}\) y la solución está elevada \(1.00\,\text{m}\) por encima del punto de inyección. Además, asuma que la viscosidad y la densidad de la solución son las del agua pura, y que la presión dentro de la vena es atmosférica.
    \vspace*{7cm}
\end{itemize}

\subsection*{SECCIÓN 10.8 Otras aplicaciones de la dinámica de fluidos}

\begin{itemize}
    \item[\textbf{30.}] Un avión tiene una masa de \(1.60 \times 10^4\,\text{kg}\) y cada ala tiene un área de \(40.0\,\text{m}^2\). Durante vuelo a nivel, la presión sobre la superficie inferior del ala es \(7.00 \times 10^4\,\text{Pa}\).
    \begin{enumerate}
        \item Suponga que el empuje sobre el avión solo se debió a una diferencia de presión. Determine la presión sobre la superficie superior del ala.
        \item Más real, una parte significativa del empuje se debe a la deflexión del aire hacia abajo causada por el ala. ¿La presión en (a) es más alta o más baja si se incluye esta fuerza? Explique.
    \end{enumerate}
    \vspace*{7cm}
    
    \item[\textbf{31.}] Un sifón se utiliza para drenar agua de un tanque, como se ilustra en la figura P10.31. Suponga flujo estable sin fricción.
    \begin{enumerate}
        \item Si \(h = 1.00\,\text{m}\), encuentre la rapidez del flujo de salida en el extremo del sifón.
        \item \textbf{¿Qué pasaría si?} ¿Cuál es la limitación en la altura de la parte superior del sifón respecto al extremo de éste? \textit{Nota:} Para que el flujo del líquido sea continuo, su presión no debe caer debajo de su presión de vapor. Suponga que el agua está a \(20.0^\circ\text{C}\), con una presión de vapor de \(2.3\,\text{kPa}\).
    \end{enumerate}
    
    % [Espacio para Figura P10.31]
    \vspace*{7cm}
\end{itemize}

\newpage

\subsection*{PROBLEMAS ADICIONALES}

\begin{itemize}
    \item[\textbf{32.}] Hace décadas se creía que los grandes dinosaurios herbívoros, como el \textit{Apatosaurus} y el \textit{Brachiosaurus}, habitualmente caminaban sobre el fondo de los lagos, y sus largos cuellos les permitían tomar aire de la superficie para respirar. El \textit{Brachiosaurus} tenía sus orificios nasales en la parte superior de su cabeza. En 1977, Knut Schmidt-Nielsen puntualizó que para tal criatura sería muy difícil el proceso respiratorio. Como un simple modelo, considere una muestra de \(10.0\,\text{L}\) de aire a presión absoluta de \(2.00\,\text{atm}\), con densidad de \(2.40\,\text{kg/m}^3\), localizada en la superficie de un lago de agua dulce. Encuentre el trabajo requerido para transportarla a una profundidad de \(10.3\,\text{m}\), permaneciendo constantes su temperatura, volumen y presión. Este requerimiento de energía es mayor que la energía que puede obtenerse mediante el metabolismo de alimentos con el oxígeno en esa cantidad de aire.
    \vspace*{7cm}
    
    \item[\textbf{33.}] Un globo (su envoltura tiene una masa de \(m_g = 0.250\,\text{kg}\)) se llena con helio y se amarra a una cuerda uniforme de longitud \(\ell = 2.00\,\text{m}\) y masa \(m = 0.050\,\text{kg}\). El globo es esférico, con un radio de \(r = 0.400\,\text{m}\). Cuando se libera en aire a temperatura de \(20^\circ\text{C}\) y densidad \(\rho_{\text{aire}} = 1.20\,\text{kg/m}^3\), se eleva una longitud \(h\) de cuerda y después permanece en equilibrio como se muestra en la figura P10.33. Se desea encontrar la longitud de cuerda levantada por el globo.
    \begin{enumerate}
        \item Cuando el globo permanece estacionario, ¿cuál es el modelo de análisis apropiado para describirlo?
        \item Escriba una ecuación de fuerzas para el globo, a partir de este modelo, en términos de la fuerza boyante \(B\), el peso \(F_g\) del globo, el peso \(F_{\text{He}}\) del helio y el peso \(F_c\) del segmento de cuerda de longitud \(h\).
        \item Realice una apropiada sustitución para cada una de esas fuerzas y resuelva simbólicamente para la masa \(m_c\) del segmento de cuerda de longitud \(h\) en términos de \(m_g\), \(r\), \(\rho_{\text{aire}}\) y la densidad del helio \(\rho_{\text{He}}\).
        \item Obtenga el valor numérico de la masa \(m_c\).
        \item Determine la longitud \(h\) numéricamente.
    \end{enumerate}
    
    % [Espacio para Figura P10.33]
    \vspace*{8cm}
    
    \item[\textbf{34.}] El peso verdadero de un objeto se puede medir en un vacío, donde las fuerzas de flotación están ausentes. Un objeto de volumen \(V\) se pesa en el aire sobre una balanza de brazos iguales, con el uso de contrapesos de densidad \(\rho\). Al representar la densidad del aire como \(\rho_{\text{aire}}\) y la lectura de la balanza como \(F_g'\), demuestre que el verdadero peso \(F_g\) es
    \[
    F_g = F_g' + \left( V - \frac{F_g'}{\rho g} \right) \rho_{\text{aire}} g
    \]
    \vspace*{7cm}
\end{itemize}

\newpage

\begin{itemize}
    \item[\textbf{35.}] A un orden de magnitud, ¿cuántos globos llenos de helio se requerirán para levantarlo a usted? Ya que el helio es un recurso irremplazable, desarrolle una respuesta teórica en lugar de una respuesta experimental. En su solución, establezca las cantidades físicas que consideró como datos y los valores que midió o estimó para ellas.
    \vspace*{6cm}
    
    \item[\textbf{36.}] \textbf{Problema de repaso.} Suponga que cierto líquido, con densidad de \(1230\,\text{kg/m}^3\), no ejerce fuerza de fricción sobre objetos esféricos. Una bola de \(2.10\,\text{kg}\) de masa y \(9.00\,\text{cm}\) de radio se deja caer desde el reposo en un tanque profundo de este líquido desde una altura de \(3.30\,\text{m}\) sobre la superficie.
    \begin{enumerate}
        \item Encuentre la rapidez con que entra la bola al líquido.
        \item Evalúe las magnitudes de las dos fuerzas que se ejercen sobre la bola mientras se mueve a través del líquido.
        \item Explique por qué la bola se mueve hacia abajo solo una distancia limitada en el líquido y calcule esta distancia.
        \item ¿Con qué rapidez la bola aparece fuera del líquido?
        \item ¿De qué modo se compara el intervalo de tiempo \(\Delta t_{\text{abajo}}\) durante el cual la bola se mueve desde la superficie hasta su punto más bajo, con el intervalo de tiempo \(\Delta t_{\text{arriba}}\) para el viaje de regreso entre los mismos puntos?
        \item \textbf{¿Qué pasaría si?} Ahora modifique el modelo para suponer que el líquido ejerce una pequeña fuerza de fricción sobre la bola, opuesta en dirección a su movimiento. En este caso, ¿de qué forma se comparan los intervalos de tiempo \(\Delta t_{\text{abajo}}\) y \(\Delta t_{\text{arriba}}\)? Explique su respuesta con un argumento conceptual en lugar de un cálculo numérico.
    \end{enumerate}
    \vspace*{8cm}
    
    \item[\textbf{37.}] Evangelista Torricelli fue la primera persona en darse cuenta de que los seres humanos viven en el fondo de un océano de aire. Él conjeturó correctamente que la presión de la atmósfera es atribuible al peso del aire. La densidad del aire a \(0^\circ\text{C}\) en la superficie de la Tierra es \(1.29\,\text{kg/m}^3\). La densidad disminuye al aumentar la altitud (a medida que la atmósfera se adelgaza). Por otra parte, si se supone que la densidad es constante a \(1.29\,\text{kg/m}^3\) hasta cierta altura \(h\) y es cero sobre dicha altura, en tal caso \(h\) representaría la profundidad del océano de aire.
    \begin{enumerate}
        \item Use este modelo para determinar el valor de \(h\) que da una presión de \(1.00\,\text{atm}\) en la superficie de la Tierra.
        \item ¿La cima del monte Everest se elevaría sobre la superficie de tal atmósfera?
    \end{enumerate}
    \vspace*{6cm}
\end{itemize}

\newpage

\begin{itemize}
    \item[\textbf{38.}] Un parámetro común que puede utilizarse para predecir turbulencia en el flujo de un fluido es el llamado número de Reynolds. El número de Reynolds para el flujo de un fluido en una tubería es una cantidad adimensional definida por
    \[
    Re = \frac{\rho v d}{\eta}
    \]
    donde \(\rho\) es la densidad del fluido, \(v\) es su rapidez, \(d\) es el diámetro interno de la tubería y \(\eta\) es la viscosidad del fluido. La viscosidad es una medida de la resistencia interna de un líquido al flujo y tiene unidades de \(\text{Pa}\cdot\text{s}\). El criterio para el tipo de flujo es como sigue:
    \begin{itemize}
        \item Si \(Re < 2300\), el flujo es laminar.
        \item Si \(2300 < Re < 4000\), el flujo está en una región de transición entre laminar y turbulento.
        \item Si \(Re > 4000\), el flujo es turbulento.
    \end{itemize}
    \begin{enumerate}
        \item Sangre de densidad \(1.06 \times 10^3\,\text{kg/m}^3\) y viscosidad \(3.00 \times 10^{-3}\,\text{Pa}\cdot\text{s}\) se modela como un líquido puro, es decir, se ignora la presencia de los glóbulos rojos. Suponga que la sangre fluye en una gran arteria de radio \(1.50\,\text{cm}\) con una rapidez de \(0.0670\,\text{m/s}\). Demuestre que el flujo es laminar.
        \item Imagine que la arteria finaliza en un solo capilar, así que el radio de la arteria se reduce a un valor mucho más pequeño. ¿Cuál es el radio del capilar que provocaría un flujo turbulento?
        \item Los capilares reales tienen radios de aproximadamente \(5\text{–}10\,\mu\text{m}\), mucho más pequeños que el valor en el inciso (b). ¿Por qué el flujo en los capilares reales no se hace turbulento?
    \end{enumerate}
    \vspace*{8cm}
    
    \item[\textbf{39.}] En 1983, Estados Unidos comenzó a acuñar la moneda de un centavo a partir de zinc revestido de cobre en lugar de cobre puro. La masa del antiguo penique de cobre es \(3.083\,\text{g}\) y el del nuevo centavo es de \(2.517\,\text{g}\). La densidad del cobre es \(8.960\,\text{g/cm}^3\) y la del zinc es \(7.133\,\text{g/cm}^3\). Las monedas nueva y antigua tienen el mismo volumen. Calcule el porcentaje de zinc (por volumen) en el nuevo centavo.
    \vspace*{6cm}
    
    \item[\textbf{40.}] \textbf{Problema de repaso.} Con referencia al dique estudiado en el ejemplo 10.4 y mostrado en la figura 10.5,
    \begin{enumerate}
        \item demuestre que el momento de torsión total ejercido por el agua detrás del dique respecto a un eje horizontal a través de \(O\) es \(\frac{1}{6} \rho g w H^3\).
        \item demuestre que la línea de acción efectiva de la fuerza total ejercida por el agua está a una distancia \(H/3\) arriba de \(O\).
    \end{enumerate}
    \vspace*{7cm}
\end{itemize}

\newpage

\begin{itemize}
    \item[\textbf{41.}] El termómetro espíritu en vidrio, inventado en Florencia, Italia, alrededor de 1654, consiste en un tubo de líquido (el espíritu) que contiene un cierto número de esferas de vidrio sumergidas con masas ligeramente distintas (figura P10.41). A temperaturas suficientemente bajas, todas las esferas flotan, pero conforme se eleva la temperatura, las esferas se hunden una tras otra. El dispositivo es un burdo pero interesante instrumento para medir la temperatura. Suponga que el tubo está lleno con alcohol etílico, cuya densidad es \(0.789 45\,\text{g/cm}^3\) a \(20.0^\circ\text{C}\) y disminuye a \(0.780 97\,\text{g/cm}^3\) a \(30.0^\circ\text{C}\).
    \begin{enumerate}
        \item Suponga que una de las esferas tiene un radio de \(1.000\,\text{cm}\) y está en equilibrio a la mitad del tubo a \(20.0^\circ\text{C}\), determine su masa.
        \item Cuando la temperatura se incrementa a \(30.0^\circ\text{C}\), ¿qué masa debe tener una segunda esfera del mismo radio para equilibrar en el punto medio?
        \item A \(30.0^\circ\text{C}\), la primera esfera ha caído al fondo del tubo. ¿Qué fuerza hacia arriba ejerce el fondo del tubo sobre esta esfera?
    \end{enumerate}
    
    % [Espacio para Figura P10.41]
    \vspace*{8cm}
    
    \item[\textbf{42.}] Una mujer drena el agua de su pecera mediante un sifón que va hacia un drenaje exterior, como se muestra en la figura P10.42. El tanque rectangular tiene área de planta \(A\) y profundidad \(h\). El drenaje se ubica a una distancia \(d\) bajo la superficie del agua en el tanque, donde \(d \gg h\). El área de sección transversal del tubo de sifón es \(A'\). Modele el agua como flujo sin fricción. Demuestre que el intervalo de tiempo requerido para vaciar el tanque es
    \[
    \Delta t = \frac{Ah}{A' \sqrt{2gd}}
    \]
    
    % [Espacio para Figura P10.42]
    \vspace*{7cm}
    
    \item[\textbf{43.}] \textbf{Problema de repaso.} Usted y su padre están diseñando una cascada para la piscina de su patio trasero. En la parte superior de la cascada hay un tanque que contiene agua que se mantiene a una profundidad \(d = 0.280\,\text{m}\) por una bomba. Como se muestra en la figura P10.43, hay una pequeña escotilla de altura \(h = 0.100\,\text{m}\) y ancho \(w = 0.150\,\text{m}\), con bisagras en la parte superior, que se puede usar para abrir y cerrar el suministro de agua a la cascada. Desea conectar un pestillo simple en el centro de la parte inferior de la escotilla y está tratando de decidir qué tipo de pestillo comprar en su ferretería local. Para tomar esa decisión, debe determinar la fuerza que el pestillo debe poder soportar para mantener cerrada la escotilla.
    \vspace*{7cm}
\end{itemize}

\newpage

\begin{itemize}
    \item[\textbf{44.}] \textbf{Problema de repaso.} Usted y su padre están diseñando una cascada para la piscina de su patio trasero. En la parte superior de la cascada hay un tanque que contiene agua que se mantiene a una profundidad \(d\) mediante una bomba. Como se muestra en la figura P10.43, hay una pequeña compuerta de altura \(h\) y ancho \(w\), con bisagras en la parte superior, que se puede usar para abrir y cerrar el suministro de agua a la cascada. Desea conectar un pestillo simple en el centro de la parte inferior de la escotilla y está tratando de decidir qué tipo de pestillo comprar en su ferretería local. Para tomar esa decisión, debe determinar la fuerza que el pestillo debe poder soportar para mantener cerrada la escotilla.
    
    % [Espacio para Figura P10.43]
    \vspace*{7cm}
    
    \item[\textbf{45.}] \textbf{Problema de repaso.} Un disco uniforme de \(10.0\,\text{kg}\) de masa y \(0.250\,\text{m}\) de radio gira a \(300\,\text{rev/min}\) en un eje de baja fricción. Se debe detener en \(1.0\,\text{min}\) mediante un freno que hace contacto con el disco a una distancia promedio de \(0.220\,\text{m}\) del eje. El coeficiente de fricción entre el freno y el disco es \(0.500\). Un pistón en un cilindro de \(5.00\,\text{cm}\) de diámetro presiona el freno contra el disco. Encuentre la presión requerida para el fluido de frenos en el cilindro.
    \vspace*{7cm}
    
    \item[\textbf{46.}] \textbf{Problema de repaso.} En una pistola de agua, un pistón impulsa el agua a través de un tubo grande de área \(A_1\) hacia un tubo más pequeño de área \(A_2\), como se muestra en la figura P10.46. El radio del tubo grande es \(1.00\,\text{cm}\) y el del tubo pequeño es \(1.00\,\text{mm}\). El tubo menor está \(3.00\,\text{cm}\) arriba del tubo mayor.
    \begin{enumerate}
        \item Si la pistola se dispara horizontalmente a una altura de \(1.50\,\text{m}\), determine el intervalo de tiempo requerido para que el agua viaje desde la boquilla hasta el suelo. Desprecie la resistencia del aire y suponga presión atmosférica de \(1.00\,\text{atm}\).
        \item Si el alcance deseado del chorro es \(8.00\,\text{m}\), ¿con qué rapidez \(v_2\) debe salir el agua por la boquilla?
        \item ¿A qué rapidez \(v_1\) debe moverse el émbolo para lograr el alcance deseado?
        \item ¿Cuál es la presión en la boquilla?
        \item Encuentre la presión requerida en el tubo mayor.
        \item Calcule la fuerza que se debe ejercer sobre el gatillo para lograr el alcance deseado. (La fuerza que se debe ejercer está arriba de la presión atmosférica.)
    \end{enumerate}
    
    % [Espacio para Figura P10.46]
    \vspace*{9cm}
\end{itemize}

\newpage

\begin{itemize}
    \item[\textbf{47.}] Un fluido incompresible, no viscoso, inicialmente en reposo en la parte vertical de la tubería que se muestra en la figura P10.47a, donde \(L = 2.00\,\text{m}\). Cuando la válvula se abre, el fluido circula a la sección horizontal de la tubería. ¿Cuál es la rapidez del fluido cuando está por completo en la sección horizontal, como se muestra en la figura P10.47b? Suponga que el área de sección transversal de toda la tubería es constante.
    
    % [Espacio para Figura P10.47]
    \vspace*{7cm}
    
    \item[\textbf{48.}] El casco de un bote experimental se levantará sobre el agua mediante un aerodeslizador montado bajo su quilla, como se muestra en la figura P10.48. El aerodeslizador tiene forma de ala de avión. Su área proyectada sobre una superficie horizontal es \(A\). Cuando el bote se remolca con una rapidez \(v\) suficientemente alta, el agua se mueve en flujo laminar de modo que su rapidez de densidad \(\rho\) promedio en la parte superior del aerodeslizador es \(n\) veces mayor que su rapidez \(v\), bajo el aerodeslizador.
    \begin{enumerate}
        \item Si ignora la fuerza de flotación, demuestre que la fuerza de sustentación hacia arriba que el agua ejerce sobre el aerodeslizador tiene una magnitud
        \[
        F = \frac{1}{2} (n^2 - 1) \rho v^2 A
        \]
        \item El bote tiene masa \(M\). Demuestre que la rapidez de despegue es
        \[
        v = \sqrt{\frac{2Mg}{(n^2 - 1)A\rho}}
        \]
    \end{enumerate}
    
    % [Espacio para Figura P10.48]
    \vspace*{8cm}
\end{itemize}

\newpage

\subsection*{PROBLEMAS DE DESAFÍO}

\begin{itemize}
    \item[\textbf{49.}] Demuestre que la variación de presión atmosférica con la altitud está dada por \(P = P_0 e^{-\alpha y}\), donde \(\alpha = \rho_0 g / P_0\), \(P_0\) es la presión atmosférica en algún nivel de referencia \(y = 0\) y \(\rho_0\) es la densidad atmosférica a este nivel. Suponga que la disminución en presión atmosférica sobre un cambio infinitesimal en altura (de modo que la densidad es aproximadamente uniforme) puede expresarse, de la ecuación 10.4, como \(dP = -\rho g dy\). También suponga que la densidad de aire es proporcional a la presión, lo cual es equivalente a suponer que la temperatura del aire es la misma para toda altitud.
    \vspace*{8cm}
    
    \item[\textbf{50.}] \textbf{¿Por qué no es posible la siguiente situación?} Una barcaza transporta, a lo largo de un río, una carga de pequeños pedazos de hierro apilados. La pila de hierro tiene la forma de un cono cuyo radio \(r\) de la base es igual a su altura central \(h\). La barcaza tiene forma cuadrada, con lados verticales de longitud \(2r\), así que la pila de hierro llena justo los bordes de la barcaza. El navío se aproxima a un puente de baja altura, y el capitán observa que la parte superior de la pila no logrará pasar por debajo del puente. El capitán ordena a la tripulación que tire al agua pedazos de hierro para así disminuir la altura de la pila. Conforme el hierro se va tirando, la pila siempre mantiene su forma cónica cuyo diámetro es igual a la longitud lateral de la barcaza. Después de un cierto volumen de hierro lanzado al agua, la nave pasa bajo el puente sin que la parte superior de la pila pegue en el puente.
    \vspace*{8cm}
\end{itemize}

\end{document}
